\chapter{Introduction and approach}
\label{ch:sq2:intro}

This document answers the supplementary question on energy gentrification issued to Protocol Policy Lab following the hearing of 1 May 2026, with particular reference to Western Sydney, where data centre clustering is most concentrated. The inquiry to which it contributes is itself a novel exercise, described by its proponents as a process to \enquote{establish a parliamentary inquiry into data centre developments in NSW, the first of its kind in Australia} \autocite{greensNSWDataCentreInquiry2026}. That novelty matters, because it means New South Wales is forming its evidentiary and regulatory baseline at the same moment that the sector is expanding fastest.

\section{Evidence base and method}
\label{sec:sq2:method}

The answer draws on two bodies of evidence. The first is the peer-reviewed and technical literature on data centre energy, water and grid integration. The second is the primary and grey literature specific to the Australian National Electricity Market and to New South Wales, including the Australian Energy Market Operator forecasts as relayed in the New South Wales Government's own consultation paper, transmission planning by Transgrid, and analysis by Energy Consumers Australia and the United States Studies Centre. Quotations throughout are reproduced verbatim from the cited source. Measured quantities are distinguished from projections, and projections are attributed to their modelling scenario wherever the source specifies one. Where no New South Wales figure exists, this is stated rather than inferred from overseas data.

\section{The standard of inference adopted}
\label{sec:sq2:inference}

This submission's conclusion rests on inference rather than direct local measurement, and the standard applied is stated plainly so that the Committee can weigh it appropriately. The argument is deliberately bounded. Protocol does not assert that energy gentrification is occurring in New South Wales, nor that it will occur. The claim is the weaker and more defensible one that the conditions which preceded the effect in comparable jurisdictions are present in Western Sydney, that the outcome is therefore a material and foreseeable risk, and that because the outcome is contingent on policy settings it is preventable. This framing is adopted precisely because the contingency strengthens rather than weakens the case for early action, since a preventable harm warrants pre-emptive safeguards while it remains preventable.
